\begin{abstract}
Se estudiaron experimentalmente los principios de Arquímedes, continuidad y Torricelli mediante dos montajes de laboratorio. En la primera parte, se sumergió progresivamente una varilla en agua y se registraron, para diez longitudes sumergidas, la lectura de una balanza de brazo triple y su diferencia respecto a una lectura de referencia. La relación lineal entre ambas variables permitió determinar una densidad experimental del agua de $\SI{0.991}{\gram\per\cubic\centi\meter}\pm\SI{0.012}{\gram\per\cubic\centi\meter}$, con un error porcentual de $0.86\%$ y un coeficiente de determinación de $R^2=0.9930$. Se consideraron incertidumbres de $\pm\SI{0.05}{\gram}$ para la balanza, $\pm\SI{0.10}{\gram}$ para la diferencia entre lecturas y $\pm\SI{0.05}{\centi\meter}$ para las mediciones de longitud y diámetro. En la segunda parte, se midieron los tiempos de vaciado de un recipiente de Torricelli para alturas iniciales de $\SI{10}{\centi\meter}$, $\SI{8}{\centi\meter}$, $\SI{6}{\centi\meter}$, $\SI{4}{\centi\meter}$ y $\SI{2}{\centi\meter}$. El ajuste experimental fue $t=88.72h^{0.651}$, con $R^2=0.9806$, y los errores porcentuales respecto al modelo teórico variaron entre $5.73\%$ y $22.48\%$, con un promedio de $14.65\%$. Los resultados respaldaron la proporcionalidad entre el volumen desplazado y la inmersión, así como la dependencia sublineal del tiempo de vaciado con la altura inicial, aunque evidenciaron desviaciones respecto a las condiciones ideales.
\end{abstract}

\begin{IEEEkeywords}
palabra clave 1, palabra clave 2, palabra clave 3, palabra clave 4
\end{IEEEkeywords}
